\section{Limitations}\label{sec:limitations}

\subsection{Reuse of Rizzo beans}

\todo{Re-iterate that reuse in signal combinators is limited (we have already shown the map and why reusing would be unsafe), maybe show stop?. Anyway, we want to highlight the problem that too many instructions are inserted in signal combinators}

\todo{The order inc/dec reset/reuse insertion. In the map combinator, in the reset/reuse algorithm we insert a $reset \; s$ which forces us to then insert a $inc \; s$ instruction. This is less than ideal and the perceus paper (either 1 or 2) they observed something similar and proposed to flip the order.}


\subsection{Later values are computed more than once per time step}

Later values are just functions, and thus, if more than one signal depend on the same Later value, then the Later value will be evaluated more than once. This is not ideal, as the Later value may be expensive to compute and will grow the heap linearly with the number of dependencies on the Later value.
In the evaluation of the signal heaps behaviour, we mention that the signals $[4]$ and $[6]$ are caused by a Later value being evaluated 2 times, once for each dependence on the same Later value.
A workaround that currently works for \texttt{Later values} is by temporarily making the later into a signal. We can do this by giving the Later value a dummy head and using the tail constructor to get the now cached Later value back:

\begin{verbatim}
    tail ("" :: mk_sig_of_channel console)
\end{verbatim}

A proper fix could require us to change the Later type to contain 2 additional fields. One field for the cached value and one for a flag to indicate which time step the cached value is for. The evaluation of the Later value would then first check if the cached value is valid for the current time step and return it if it is, and otherwise compute the Later value, store it in the cache, and return it.
